\documentclass[11pt,letterpaper]{article}

%% ---- Packages ----
\usepackage[utf8]{inputenc}
\usepackage[T1]{fontenc}
\usepackage{mathpazo}
%% Matches the manuscript. Every table here is placed with [H], so a table box that does not
%% fit the space left on a page pushes a break and leaves the rest of that page blank; the
%% extra page height lets two of the shorter boxes share a page instead.
\usepackage[left=2cm,right=2cm,top=1.3cm,bottom=1.3cm]{geometry}
\usepackage{graphicx}
\usepackage{float}
\usepackage{caption}
%% labelfont=bf + labelsep=period make the AUTO-generated label render as
%% "\textbf{Extended Data Fig. 7.}", which is byte-for-byte what the hand-written prefix inside
%% each \caption used to produce. The prefix has been removed from all seven captions; keeping
%% both printed "Extended Data Fig. 7: Extended Data Fig. 7." on every figure. Numbering is now
%% the counter's, so a caption can no longer disagree with its own \label.
\captionsetup{singlelinecheck=false, font={small}, skip=4pt, labelfont=bf, labelsep=period}
%% Compaction, matching the manuscript's figure floats. The blank tails on the table pages come
%% from the [H] boxes, not from the tables themselves, so the space is recovered from the glue
%% around them and from the row stretch rather than by shortening any table.
\setlength{\textfloatsep}{8pt plus 2pt minus 2pt}
\setlength{\intextsep}{8pt plus 2pt minus 2pt}
\setlength{\floatsep}{8pt plus 2pt minus 2pt}
\setlength{\abovecaptionskip}{4pt}
\setlength{\belowcaptionskip}{0pt}
\renewcommand{\arraystretch}{1.15}
%% Supplementary-table blocks. Heading, tabular and note are set INSIDE one non-floating [H]
%% table so the three move as a single box. Set as running text they did not: an [H] table that
%% does not fit the remaining page is carried over whole while the heading above it stays
%% behind, which is how the headings of Supplementary Tables 1, 3 and~4 came to print a full
%% page above the tables they name. \parbox keeps the surrounding \centering from
%% ragged-centring the prose.
\newcommand{\stabhead}[1]{{\noindent\large\bfseries #1\par}\medskip}
\newcommand{\stabsub}[1]{{\noindent\bfseries #1\par}\smallskip}
\newcommand{\stabbody}[1]{\noindent #1\par\medskip}
\newcommand{\stabnote}[1]{\par\medskip\noindent{\footnotesize #1\par}}
\usepackage{amsmath,amssymb}
\usepackage{booktabs}
%% Supplementary tables are longtables. Each used to be a non-floating [H] box 380 to 430 pt tall,
%% so a table that did not fit the space left on a page was carried over whole and left the rest
%% of that page blank; six of the eight table pages were ending 26 to 46%% short. A longtable
%% splits at a row boundary instead. The heading travels inside \endfirsthead and a continuation
%% line inside \endhead, which is what keeps a heading from printing a page above the table it
%% names. Each heading element is its own \caption* row: \caption* is not \long, so one caption
%% cannot hold the \par that a title and a sub-title together would need.
\usepackage{longtable}
\setlength{\LTcapwidth}{\textwidth}
\setlength{\LTpre}{0pt}
\setlength{\LTpost}{\medskipamount}
\usepackage[dvipsnames]{xcolor}
\usepackage[most]{tcolorbox}
\usepackage{hyperref}
\usepackage[numbers,sort&compress,super]{natbib}
\usepackage{enumitem}
\usepackage{lineno}
\linenumbers
%% Same value as the main manuscript (seed-resampling control range).
%% Macros mirrored from the main manuscript for the Supplementary Notes below.
%% Values must stay identical to PopRetrieve_manuscript.tex.
%%
%% MAIN-FIGURE REFERENCES ARE HARD-CODED HERE, AND HAVE TO BE. This is a separate document,
%% so \ref{fig:N} does not resolve against the manuscript's labels. Every reference this
%% file makes into the main figures therefore goes through one of the two macros below,
%% which is the single place to edit if the main deck is ever renumbered or relettered.
%% Checked against PopRetrieve_manuscript.tex on 2026-09-01.
%% THIS ONE WENT STALE AND IS THE REASON THE BLOCK ABOVE EXISTS. It read "Fig.~5e" until
%% 2026-09-04. Figure 5e is now the headroom-correlation null, Figure 5f is the recoverability of
%% the Tahoe state that carries the drug-by-state interaction, and the CONSTRUCTED two-state
%% mixture this note is about has no figure panel at all: it is reported in the main Results as
%% prose. A figure number is therefore the wrong kind of pointer for it, so the macro names the
%% result by its two values instead, which cannot be renumbered.
\newcommand{\MAINRECOVER}{ceiling 0.692 against best unsupervised 0.674}
\newcommand{\MAINCLASSAB}{Fig.~3a,b} % the 600-query response-matching / MoA pair
\newcommand{\NCQUERIES}{103}
\newcommand{\PREMISEDISJOINT}{+0.835}
\newcommand{\TAHOEBELOWMIX}{0.11}
\newcommand{\TAHOECONDS}{3630}
\newcommand{\TAHOECOUPLING}{0.620}
\newcommand{\TAHOECOUPLINGN}{44}
\newcommand{\TAHOEGMM}{0.582}
\newcommand{\TAHOEGONE}{0.739}
\newcommand{\TAHOEGONEQ}{0.600 to 0.836}
\newcommand{\TAHOEGTHREECC}{0.841}
\newcommand{\TAHOEGTHREECCN}{44}
\newcommand{\TAHOEGTHREEST}{0.778}
\newcommand{\TAHOEGTHREESTMIN}{0.236}
\newcommand{\TAHOEGTHREESTN}{45}
\newcommand{\TAHOEGTHREESTOPEN}{2}
\newcommand{\TAHOEGTWOCEIL}{0.853}
\newcommand{\TAHOEGTWOGAP}{0.157}
\newcommand{\TAHOEGTWOLINES}{48}
\newcommand{\TAHOEGTWOLINESPOS}{48}
\newcommand{\TAHOEGTWOPAIRS}{960}
\newcommand{\TAHOEGTWOUNSUP}{0.611}
\newcommand{\TAHOEKMEANS}{0.592}
\newcommand{\TAHOELEIDEN}{0.563}
\newcommand{\TAHOEOPENGONE}{0.394 to 0.534}
\newcommand{\TAHOEOPENGTHREE}{0.605 to 0.764}
\newcommand{\TAHOEPARTAGREE}{0.466}
\newcommand{\CONTROLRANGE}{1 to 4}
\newcommand{\TODO}[1]{{\color{BrickRed}\textbf{[[#1]]}}}

%% ---- Section formatting (Nature Methods style) ----
\usepackage{titlesec}
\titleformat{\section}{\large\bfseries}{}{0em}{}
\titleformat{\subsection}{\normalsize\bfseries}{}{0em}{}
\titlespacing*{\section}{0pt}{1.5ex plus 0.5ex}{0.8ex plus 0.2ex}
\titlespacing*{\subsection}{0pt}{1.2ex plus 0.3ex}{0.5ex plus 0.1ex}

\hypersetup{colorlinks=true,linkcolor=MidnightBlue,citecolor=MidnightBlue,urlcolor=MidnightBlue}
\setlength{\emergencystretch}{3em}

%% ONE TITLE, TWO DOCUMENTS. This must be the main manuscript's title verbatim. It was not:
%% it carried "Single-cell population structure adds limited decision value beyond response
%% magnitude in drug retrieval", which the main text stopped using when the title became a
%% question. Checked against PopRetrieve_manuscript.tex on 2026-09-04.
\title{Supplementary Information for\\
\textit{When do single-cell response distributions improve drug retrieval?}}

\author{}
\date{}

\renewcommand{\figurename}{Extended Data Fig.}

\begin{document}
\maketitle

\noindent
This Supplementary Information provides additional methodological detail, robustness analyses and
supporting experiments for the main text. It contains Supplementary Methods, four Supplementary
Notes and four Supplementary Tables, and no Extended Data figures. The supplementary analyses follow the same
conceptual structure as the main text: they examine the information retained by population-level
retrieval, the dependence of apparent retrieval gains on the evaluation criterion, and the
conditions under which state-specific single-cell responses can become consequential for candidate
ranking.

\bigskip








%% =====================================================================
%%  SUPPLEMENTARY NOTES
%%  One barrier here, and only here: it flushes every Extended Data float so
%%  the Notes begin after the figures rather than wrapping around them.
%% =====================================================================
\clearpage

\noindent{\large\textbf{Supplementary Methods}}
\bigskip

\noindent
\textbf{Implementation of population retrieval scores.}
Energy distance was the primary population score. RBF-kernel maximum mean discrepancy used the
median squared pairwise distance of the pooled query and candidate sample as its bandwidth
reference and evaluated scales of 0.25, 1 and 4 times that value. Sliced Wasserstein-1 distance
averaged quantile-matched projections over 128 fixed random directions. Candidate populations were
subsampled with a fixed seed to at most 500 cells in the observed-population analyses and 300 cells
in the partial-observation analyses. A population containing fewer than two cells after filtering
was assigned a fixed large penalty distance rather than removed from the candidate library.

\noindent
\textbf{Implementation of subpopulation coverage.}
Query cells were partitioned by $k$-means with $k=2$, three initializations and a fixed seed.
Candidate cells were assigned to the query component whose centroid had the highest cosine
similarity. A candidate with no cell assigned to a query component received the penalty distance
for that component. At small $\beta$, the coverage distance was evaluated using
$D_\beta=\operatorname{mean}(d)+\tfrac{\beta}{2}\operatorname{var}(d)+O(\beta^2)$ to avoid
numerical cancellation in the log-sum-exp form. Coverage-mean and coverage-worst use the
$\beta\rightarrow0$ and $\beta\rightarrow\infty$ limits, respectively. The base distance was
energy distance unless otherwise specified; the complete score registry is given in
Supplementary Table~1.

\bigskip

\noindent{\large\textbf{Supplementary Note 1. Additional diagnostics of when population information changes ranking}}
\bigskip

\noindent
The main text separates the existence of state-specific perturbation responses from their
recoverability and from the downstream requirement that they alter candidate ordering. We
performed several additional analyses to ask whether simple observables could identify queries
for which population-level retrieval is more useful.

\paragraph{Query identifiability.}
In A549 leave-drug-out retrieval, query identifiability was quantified as a two-cluster
silhouette on the query cells, independently of the retrieval outcome. Across 543 evaluable
query--seed combinations it was weakly and non-monotonically associated with the absolute
mechanism-of-action nDCG of a single energy-distance run (Spearman $\rho=0.17$,
$p=6.3\times10^{-5}$), and the association is confounded with candidate-library composition: the
highest-identifiability tertile holds a median of six same-mechanism candidates against two in each
lower tertile, and a larger relevant pool mechanically raises the attainable nDCG. The median
silhouette score even in that tertile is 0.038, so the two-state structure these queries carry is
weakly resolved. This analysis scores one retriever and therefore bears on absolute recovery rather
than on the population-minus-mean gain.

\paragraph{Unsupervised recovery of constructed two-state structure.}
Response-identity recoverability in the constructed two-state mixture is a supervised ceiling
against the best unsupervised method (\MAINRECOVER), and separately a ceiling of 0.879 when the
same cells are posed at the level of individual drug pairs rather than pooled drug classes. Three further
measurements characterise the same limit.

First, across nine unsupervised clustering configurations, the median adjusted Rand index against
the known two-source labels ranged from 0.037 to 0.106. The configurations were $k$-means with
$k=2$ in the raw, 10-dimensional PCA, 50-dimensional PCA and response spaces, two-component
Gaussian mixtures in the raw and 50-dimensional PCA spaces, and $k$-means with $k$ selected from
3, 4 and 5 by silhouette, reported separately by which $k$ was selected. The highest value,
0.106, was obtained when silhouette selection chose $k=4$. Recovery is weak under all nine, so the
low ceiling is not specific to the clustering method used in the main analysis.

Second, the weak separation is visible in the raw geometry rather than only in an accuracy score.
Silhouette scores computed independently for each perturbation-response population under a
two-cluster partition had a median of 0.034 in SciPlex3 ($n=529$ populations) and 0.084 in
Frangieh ($n=231$); in both datasets fewer than 0.5\% of populations exceeded 0.1. Visible
multimodality is therefore not a prerequisite for, and should not be equated with, recoverable
state-specific perturbation information.

Third, recovery is limited by state separation rather than by cell number. Averaging the
$k$-means $k=2$ adjusted Rand index over 20 seeds across a grid of source separation and cells per
source, increasing the cell budget from 25 to 400 cells per source changed recovery by at most
0.02 at fixed separation, whereas increasing separation from 1.0 to 8.0 times the observed value
raised it from 0.07 to 1.00. The separation transform was defined from the full source
populations before subsampling. Additional cells therefore do not compensate for states that are
close together, which is the same conclusion the supervised ceiling reaches by a different route.

%% REISSUED 2026-09-04 under the U arm, from
%% results/exp16_gate_diagnosis/_merged_query_divergence.csv, which was itself rebuilt on
%% 2026-09-03 (see docs/phase2/POST_REPAIR_MASTER_RESULTS.md A5c). The V-arm values were
%% rho = 0.14, permutation p = 2e-4, a top-stratum mean of about +0.02, and a reversal in one of
%% three cell lines. The reversal is now in TWO of three, which is a change of claim and not only
%% of digits, so the cell-line means are printed rather than summarised.
\paragraph{Measured response divergence.}
We next stratified the 600 mechanism-evaluable queries by measured response divergence, one minus
the mean cosine between their state-specific perturbation responses. Greater divergence was
associated weakly with a larger population-minus-mean mechanism-recovery difference
(Spearman $\rho=0.15$, $p=1.4\times10^{-4}$), but the conditional
advantage did not become convincingly positive. The most divergent quartile reached a mean nDCG
difference of $+0.014$ with a bootstrap interval of $-0.012$ to $+0.041$, and the direction
reversed in two of the three SciPlex3 cell lines (A549 $-0.107$, $n=9$; K562 $-0.021$, $n=38$;
MCF7 $+0.040$, $n=96$). That quartile is not balanced across cell lines, and its positive mean is
carried by the one line that contributes two thirds of it.

These results distinguish differential response from decision relevance. Stronger state-specific
responses may create the information required for population-aware retrieval, but they do not
guarantee that this information changes the ordering of candidate perturbations.

\paragraph{Precision of the two task-proximal comparisons in the highest-divergence stratum.}
In the most divergent quartile the two criteria differ in precision rather than in effect size.
The mechanism-recovery gap is nearly three times the minority-coverage gap but its per-query
standard deviation is twelve times larger ($0.169$ against $0.0136$), so the two readings are
asymmetric: the minority-state comparison is well determined and returns a small effect, which is
evidence of a small effect, whereas the mechanism-recovery comparison in this stratum supports
uncertainty rather than evidence of no effect. Achieved power and the sample size required for
80\% power are given in the code release. The main text's mechanism-recovery conclusion rests on
the full 600-query set and on the shape of the whole distribution (\MAINCLASSAB), not on this
stratum.

\paragraph{Instability of thresholded dominance counts.}
Counting tasks whose population-level gain exceeds a fixed threshold is unstable under resampling:
at a gain threshold of $+0.01$, only 2 of the 10 initially threshold-crossing cross-line tasks stay
above it in at least 8 of 10 seed resamples, while control drugs that cross it in no recorded seed
cross it in \CONTROLRANGE\ of 10 resamples. No count of dominant tasks is therefore reported from
this work, and the main text reports continuous retrieval differences and stratified analyses
instead.


\bigskip
%% REISSUED 2026-09-04 UNDER THE UNBIASED U-STATISTIC ENERGY DISTANCE. Every number in this
%% paragraph was computed under the V-statistic and none of them survived the repair: the counts
%% were 621 / 133 / 11 and are now 627 / 127 / 11, because the gate's verdict is computed from
%% quantities that pass through the energy kernel and six queries crossed its threshold. The
%% main manuscript reissued its copies on 2026-09-03 (see \NREC, \NNONREC there) and this
%% paragraph did not, which is what left the two documents quoting different denominators for
%% the same 765 queries. Recomputed here from
%% results/exp12_partial_observed_retrieval/per_query_scores.csv with the estimator, seed and
%% n_boot of figures/fig3/fig3e.py, so the SI, the main text and the panel share one bootstrap.
\paragraph{Two checks on the diagnostic's decline verdicts.}
The pre-specified diagnostic declines 127 queries as ``mean or no call'' and a further 11 as
``mean sufficient''. Pooling the 11 into the non-recommended group leaves the verdict unchanged,
and every interval involved is wider than either group median, so the comparison is a failure to
detect enrichment rather than a demonstration that there is none. The two decline verdicts behave
oppositely: the 127 ``mean or no call'' queries are \emph{more} divergent than the queries the
diagnostic recommends (medians 1.68 against 1.66), which is the direction its design does not
predict, while the 11 ``mean sufficient'' queries lie below the fifth percentile of both groups,
where the diagnostic is correct.

\bigskip
\noindent{\large\textbf{Supplementary Note 2. Verification and transfer limits of HIR-Bench}}
\bigskip

\noindent
HIR-Bench was introduced as a controlled construction in which the conditions under which
within-population structure can alter candidate preference are known analytically. It is therefore
used to verify the logic of a preference-changing regime rather than to infer a biological phase
boundary from data.

The generator satisfies its internal consistency checks on the full parameter grid. All three
pre-specified checks pass: when the two subpopulations have no conflict in candidate preference the
retrieval-flip rate is 0 against a threshold of 0.05, as required by construction; the median
distance between the grid's own mixing fraction and the analytical boundary
\[
\alpha^{*}
=
\frac{B}{A+B},
\]
is 0.292, a range guard rather than a discriminating test since the statistic cannot exceed its
threshold of 1.0; and under a predicted-mean information condition the
population-level and mean-level Hit@1 differ by 0.051 against a threshold of 0.10, which is the
degeneracy the construction requires when no within-population structure is available to the
retriever.

The measured counterpart of the analytical boundary behaves as the construction predicts. Over the
36 cells of the conflict-by-majority-fraction grid, the energy distance exceeds mean-signature
retrieval in Hit@1 in 30 of them, and the advantage runs from $-0.25$ to $+1.00$, largest where
conflict is high and the majority fraction is above $\alpha^{*}$ and vanishing below it.

We next asked whether the synthetic regime could be recognized from quantities observable at
retrieval time. Observable features included the subpopulation-variance ratio, isotropy index and
response diversity of the query and candidate populations. Under a cross-validation unit that
respects the 28 label-determining parameter cells, those features reach a pooled out-of-fold AUC
of 0.788 against a majority-class rate of 0.643. Splitting instead on the 672 grid configurations,
which leaves about twelve near-duplicates of each held-out instance in training, raises it to
0.836, so leakage alone is worth $+0.048$ and the honest observable result is modest. Because the
label is deterministic within a parameter cell, every held-out fold is single-class and the pooled
out-of-fold AUC is the only statistic the design admits. Removing within-population information
also raised the best decision regret over the method set, which a predicted-structure condition
recovered, so the construction penalizes that loss in the direction it was built to.

The analytical boundary was also projected to the biological retrieval tasks as a transfer test.
This projection did not generalize. The resulting labels were strongly imbalanced
(229 versus 10), for which a constant majority-class predictor achieves 95.8\% accuracy. The
synthetic-regime projection reached 39.7\%, far below this baseline and below it on four of the
five datasets separately (0.29, 0.40, 0.33 and 0.22); the fifth is the predicted-mean subset, where
the projection returns the same verdict for every task by construction. The boundary was transferred by percentile rather than by value: the
crossover percentile was fitted on 4{,}480 grid cells of the HIR-Bench method-performance layer
and applied to the real tasks at the corresponding quantile of their own preference-conflict
distribution.

HIR-Bench establishes constructively that within-population structure can become consequential
for candidate choice. The failed transfer sets its limit: the corresponding regime in biological
data cannot be inferred from this synthetic boundary.




\bigskip
\noindent{\large\textbf{Supplementary Note 3. Robustness of the natural-heterogeneity and large-scale analyses}}
\bigskip

\noindent
The main text uses patient-derived glioblastoma and Tahoe-100M to ask whether the information
requirements identified in the controlled experiments are also present when cellular heterogeneity
is not created by combining experimentally distinct source populations. These datasets measure differential response,
recoverability and the extent to which state-specific response ordering leaves headroom for
population structure to alter candidate preference; no retrieval is benchmarked on them.

\paragraph{Response alignment between compartments.}
For 17 patient--drug pairs in which both compartments carried enough treated and control cells, the
cosine between the malignant and myeloid perturbation-response directions has a median of 0.566,
against a range of 0.02 to 0.39 for the constructed mixtures used elsewhere in this study. Read
directly, that contrast says the two compartments of a real tumour respond far more alike than two
experimentally combined cell lines do.

The cosine is descriptive. It is confounded with effect size and with signal-to-noise: a
compartment with a weak or noisily estimated response yields a lower cosine for reasons that have
nothing to do with its pharmacology, and the compartments here differ substantially in cell
number. The two natural-tissue results the main text relies on, the recoverability gap and the
conservation of candidate ordering across disjoint compartments, are measured without reference
to it.

\paragraph{Robustness of glioblastoma compartment assignment.}
The ZhaoSims2021 dataset does not provide resolved cell-type labels for the analysed cells, so
malignant, tumour-associated myeloid and oligodendrocyte compartments were assigned from marker
scores. Cells with fewer than 200 detected genes were removed, counts were library-size normalized
to $10^4$ and $\log(1+x)$ transformed, and gene symbols were upper-cased; compartment scoring runs
on that full matrix, before the 2,000-gene selection used downstream, so the labels do not depend
on a gene set chosen for downstream variance. Three marker sets were used: malignant \emph{SOX2},
\emph{EGFR}, \emph{PDGFRA}, \emph{OLIG1}, \emph{OLIG2}, \emph{GFAP} and \emph{AQP4}; myeloid
\emph{PTPRC}, \emph{CD14}, \emph{AIF1}, \emph{CSF1R}, \emph{ITGAM}, \emph{C1QA} and
\emph{C1QB}; oligodendrocyte \emph{MBP}, \emph{PLP1}, \emph{MOG} and \emph{MAG}. A marker
absent from the matrix is skipped. Each gene was $z$-scored across all cells of the cohort and the
compartment score is the mean of its set's $z$-scores, so the scores are cohort-dependent by
construction. A cell was called for the highest-scoring compartment only if the mean
log-normalized expression of that compartment's key markers (malignant \emph{SOX2},
\emph{OLIG2}, \emph{EGFR}; myeloid \emph{CD14}, \emph{CSF1R}, \emph{C1QB}, \emph{AIF1};
oligodendrocyte \emph{MBP}, \emph{PLP1}) reached 0.25 and its $z$-score margin over the
second-highest compartment reached 0.25. Cells failing either criterion were left unassigned
rather than swept into a leftover class, which leaves 39.9\% of cells unassigned. The build
refuses to proceed unless every called compartment expresses its own key markers most strongly,
and writes the check to \texttt{results/zhao\_gbm/compartment\_validation.csv}.

The results were stable when both thresholds were varied jointly over
\[
\{0.10,\;0.25,\;0.40,\;0.50\}.
\]
The unsupervised arm is the best of four partitions of the same cells: $k$-means with $k=2$,
a two-component Gaussian mixture, Leiden\cite{Leiden} and HDBSCAN\cite{HDBSCAN}. Leiden is the
maximum on 152 of the 180 natural splits, so the reported best-unsupervised value is close to
Leiden's alone.

Across this range, the excluded fraction varied from 29\% to 58\%, while the supervised
drug-response ceiling remained between 0.915 and 0.923, the best unsupervised recovery between
0.77 and 0.80, and the paired recoverability gap between approximately 0.11 and 0.13. The median
cosine between malignant- and myeloid-specific perturbation-response directions remained between
0.54 and 0.59. Thus, the natural-tissue conclusions are not specific to the primary 0.25
assignment threshold.

\paragraph{Patient-level uncertainty.}
The natural-tissue recoverability analysis contains 36 within-patient drug-pair splits across four
patients and two compartments, with 30 of the 36 splits contributed by a single patient. Drug-pair
splits were therefore not treated as independent replicates. Patient-level cluster bootstrap
resampling was used to quantify uncertainty, and the analysis was repeated after removing the
dominant patient.

The reported recoverability gap is computed within each split before aggregation,
\[
G_j
=
A^{(j)}_{\mathrm{sup}}
-
A^{(j)}_{\mathrm{unsup}},
\]
rather than as the difference between the two marginal medians. The primary paired median is
$+0.117$; subtracting the marginal medians would instead give $+0.147$. Across 5,000
patient-cluster bootstrap resamples, the natural-minus-constructed difference in recoverability
gap was $+0.109$ with an interval of $[+0.052,+0.129]$, and no bootstrap draw was at or below
zero. Because only four patient clusters contribute to this calculation, the interval is
interpreted as supportive rather than as a precise population-level uncertainty estimate.

Two aggregation controls give the same gap. Fixing every probe--clusterer pair, so that the same pair is scored in both arms,
gives $+0.122$ in tissue against $-0.002$ in the constructed arm. Removing the patient that
contributes 30 of the 36 splits gives $+0.123$, with a $[+0.098,+0.138]$ interval.

\paragraph{Large-scale measurement in Tahoe-100M.}
Tahoe-100M plate 3 provides a substantially larger setting in which heterogeneity occurs within
cell lines rather than through constructed mixtures. Two cellular-state definitions are used
below, and both are fixed before any treated cell is examined. The first is the released
cell-cycle phase call, G1 against G2M, with S-phase cells dropped as intermediate and each
state's response referred to its own phase-matched control. The second is $k$-means with $k=2$
fitted on each cell line's DMSO\_TF control cells in the 2,304 highly variable genes in log1p
space, with ten initializations and a fixed seed; treated cells are assigned to the nearest
control centroid, the larger control cluster is taken as the first state, and a cell line enters
only when it holds at least 100 control cells. Across 3,630 evaluable drug--cell-line
conditions, the median cosine between the two state-specific response directions was 0.739
(interquartile range 0.600--0.836). Fewer than 0.11\% of conditions were as divergent as the
constructed mixtures used in the controlled experiments.

This estimate is not driven materially by the minimum-cell filter. Conditions were excluded when
one of the required state-specific treated or control arms contained fewer than 50 cells, which
can preferentially remove strong cell-cycle-arresting perturbations. Across drugs, coverage was
only weakly associated with induced-response cosine (Spearman $\rho=0.077$), and restricting the
analysis to well-covered compounds changed the median from 0.739 to 0.740.

Recoverability remained measurable at scale. Across 960 sampled drug pairs and 48 cell lines, the
supervised ceiling was 0.853 and the best unsupervised procedure reached 0.611, giving a median
paired gap of 0.157. The gap was positive in all 48 evaluable cell lines. Individual unsupervised
methods gave similar results: $k$-means 0.592, Gaussian mixtures 0.582 and
Leiden\cite{Leiden} 0.563.

\paragraph{State-specific ordering and decision-relevance headroom.}
We finally compared relative drug-response ordering between non-overlapping cellular states. Under
the cell-cycle partition, the median disjoint Spearman correlation between the two drug--drug
similarity structures was 0.841 across 44 evaluable contexts. Under the independently
control-derived two-state partition, the corresponding median was 0.778 across 45 contexts.

These high correlations indicate that relative drug ordering is often substantially preserved
between cellular states even when the states respond differently. The pattern is not universal:
under the control-derived partition, 2 of 45 contexts had a correlation below 0.5, with a minimum
of 0.236. Such contexts leave greater headroom for state-specific response information to alter
candidate ordering. Headroom of that kind is what a population-level method would need in order to
change the selected candidate, and it is measured here on the ordering alone.

The two state definitions agree only moderately in their ranking of contexts
(Spearman $\rho=0.466$), reinforcing that apparent decision-relevance headroom depends on how
cellular states are defined. Both partitions were specified independently of downstream retrieval
performance.

Tahoe-100M therefore supports the same separation of concepts observed elsewhere in the study:
state-specific responses can exist, and response identity can remain recoverable, while much of
the relative candidate ordering is still preserved across cellular states. These are intrinsic-state
diagnostics: they constrain the conditions under which population-level retrieval could add value.
The advantage itself is measured on the same plate by the source-to-target intervention benchmark
of the main text (Figs.~5 and~6).





\noindent{\large\textbf{Supplementary Note 4. Perturbation-predictor diagnostics}}
\bigskip

\noindent
The predictor analyses ask one question: whether a predicted response population contains
state-specific perturbation effects that a downstream population-aware retriever could exploit.

\paragraph{Additive response models.}
For a predictor that applies the same perturbation effect $\delta_d$ to every cell in a context,
\[
\widehat{x}_{i,d}
=
x_i+\delta_d,
\]
the response of any two cellular states is identical after subtraction of their corresponding
controls,
\[
r_{d,0}
=
r_{d,1}
=
\delta_d.
\]
Consequently,
\[
\cos(r_{d,0},r_{d,1})=1.
\]
This result is independent of the heterogeneity of the starting control population. An additive
predictor can therefore generate a visually heterogeneous predicted population while containing
no state-specific variation in the perturbation effect itself.

The average-effect and linear-latent predictors used in the main analysis satisfy this additive
limit by construction when the perturbation response is referenced to the model-consistent
baseline. They are included to establish this limiting case.

%% MOVED HERE FROM THE MAIN METHODS, 2026-09-04. It specified an experiment that appears only in
%% this note and in Supplementary Table 4, and while it sat in the main Methods it claimed to
%% specify Fig. 6c, which is the Tahoe-100M prediction arm and a different experiment entirely.
\paragraph{The predict-then-rank construction.}
The predictor-generated retrieval used in this note factorizes the pipeline the field actually uses
into a forward predictor and a retrieval rule, so that the retrieval layer can be varied with the
predictor held fixed. It is run on SciPlex3 and is not the Tahoe-100M prediction arm of the main
text. Queries are cross-line: for each of the three SciPlex3
cell-line pairs (K562--A549, A549--MCF7, K562--MCF7) and each of ten drugs selected by a
cross-context divergence probe, the query is that drug's response blended across the two lines at
mixing fraction $\alpha\in\{0.5,0.7,0.9\}$ over 200 cells, with eight seeds, giving 720 queries.
The candidate library is the query's own drug together with up to 30 further drugs shared by both
lines; the query's drug is the single relevant candidate. Each predictor is asked to imagine every
candidate's response in the query's real blended control cells, mixed at the same $\alpha$, so no
predictor sees the within-query line structure. The imagined populations are then ranked by four
retrieval rules: mean cosine, $R^2$ of the predicted against the observed mean delta, energy
distance and coverage. Performance is nDCG@10 with binary relevance, and the reported quantity is
the distributional minus the mean-cosine value for the same predictor, so each predictor is its own
control. The predictors are fitted on the full SciPlex3 matrix, including the lines that compose
each query, so these values are an in-sample ceiling on candidate generation rather than a transfer
result.

\paragraph{Two reference baselines.}
A predicted population's induced response can be referred to two different controls, and the two
answer different questions. Referring it to the model's own zero-effect baseline, which is its
reconstruction of the control population whenever the model decodes and the control cells
themselves when it does not, cancels the reconstruction bias and leaves the divergence the model
itself generates. The additivity result above pins this quantity to exactly 1 for any predictor
of the form $\widehat{x}_{i,d}=g(x_i)+\delta_d$ with a fixed $g$, covering both the gene-space
case $g=\mathrm{id}$ and the latent case in which $g$ is the model's own reconstruction, so the
value is analytic and not estimated. Referring it to the real
control cells is instead what a retrieval scorer does, and the resulting cosine contains the
model's reconstruction bias, which differs between cellular states and is not a perturbation
response. Across the same 30 cell-line and drug combinations, the average-effect, linear-latent,
scGen and optimal-transport predictors give 1.000, 1.000, 0.972 and 0.900 against the model
baseline, and 1.000, 0.190, 0.404 and 0.380 against the real control; real treated cells give
0.205, and have no model baseline.

The two columns are informative together. Lower cosine means more state-specific divergence. The
average-effect predictor adds its effect directly in gene space and has no reconstruction step,
so its two references coincide at 1.000, which is the correct answer for such a model. The other three pass their output through a decoder whose
reconstruction error differs between cellular states, and that error enters the real-control
cosine as if it were a perturbation effect. The linear-latent model is the clearest case: it
generates exactly zero state-specific divergence, yet measured the way a scorer measures it it
reads 0.190, marginally more divergent than the 0.205 of the real cells it replaces. A
distributional score run on such a predicted population is therefore partly scoring the decoder.
This is a second reason, independent of transfer accuracy, why predicted candidates do not
provide a clean test of whether population structure improves retrieval.

\paragraph{scGen.}
A deployed scGen model was evaluated using the same within-context cellular-state partition and
state-specific induced-response definition as the other predictors. Across the same 30
cell-line and drug combinations, its mean induced-response cosine against its own
reconstruction baseline was 0.972, close to the additive limit of 1.000, and 0.404 against the
real control. Thus, in this experiment scGen preserved substantially less state dependence in
the perturbation effect than the 0.205 observed in real treated cells, under either reference.

The model is the published \texttt{scgen.SCGEN} implementation on \texttt{scvi-tools}\cite{scviTools},
trained on
the SciPlex3 tensor with the perturbation label as batch key and the cell line as label key, 30
latent dimensions, batch size 512, at most 40 epochs with early stopping after 15 epochs without
improvement, and a fixed seed. Predictions use scGen's own latent arithmetic: one latent delta per
perturbation, taken as the mean latent of treated cells minus the mean latent of control cells,
added to the encoded control mean and decoded.

\paragraph{Entropic optimal-transport mapping.}
We additionally evaluated an entropic optimal-transport map as a predictor that is not
constrained to apply a shared additive shift to all cells. It is a different map from the
Gaussian (Bures) predictor of the Tahoe-100M arm (Methods). Under the within-context
differential-response analysis, the optimal-transport predictions reached an induced-response
cosine of approximately $0.900\pm0.107$ against their own reconstruction baseline, moving away
from the additive limit but remaining much more aligned than the observed responses.

A latent-space implementation passed the corresponding perturbation-learning diagnostic
(mean-delta agreement 0.352), whereas the same construction applied directly in gene space showed
substantially weaker perturbation learning (0.057). The latent-space implementation was therefore
used for the differential-response comparison.

The map is an entropic (Sinkhorn) coupling\cite{Sinkhorn} between control and treated cells,
solved with POT\cite{POT}, with uniform
marginals, regularization $0.05$, at most 500 iterations and a cost matrix rescaled by its maximum,
solved in a 30-component PCA latent fitted unsupervised on at most 20,000 cells. Each fitting
context is centred on its own control mean, and the query context is excluded from the fit. A
barycentric projection gives a per-source-cell displacement, extended to the query's real control
cells by nearest-source matching in the same centred space and decoded through the PCA inverse
transform.

Opening this state-specific response channel was not sufficient to provide a clean retrieval test.
When the same entropic map was used to generate candidates in the cross-context
predict-then-rank setting, Hit@1 was 0.053 under mean cosine and 0.007 and 0.008 under the two
distributional rules, against 0.879 to 0.899 for the additive average-effect candidate source. The optimal-transport predictions were
therefore too inaccurate in that transfer setting to test whether greater state dependence itself
improves retrieval.

Taken together, these analyses leave a specific unresolved experiment. A decisive test requires a
predictor that both generalizes accurately to the target cellular context and reproduces the
state-specific perturbation responses present in real cells. No predictor evaluated here satisfied
both requirements simultaneously.


\bigskip
% SUPPLEMENTARY TABLE 1
%% No \clearpage here. Supplementary Table 1(a) is a 400 pt [H] box and the Supplementary Methods
%% end near the top of their last page, so the break was throwing away most of a page that the
%% first table fits on.
\renewcommand{\thetable}{S\arabic{table}}
\setcounter{table}{0}

\begingroup
\footnotesize
\begin{longtable}{@{}p{3.2cm}p{4.0cm}p{4.3cm}p{4.0cm}@{}}
\caption*{{\large\bfseries Supplementary Table 1. Retrieval scores, task constructions and evaluation criteria.}}\\
\caption*{{\bfseries (a) Retrieval scores.}}\\
\toprule
Score & Input representation & Ranking rule & Role \\
\midrule
\endfirsthead
\caption*{\textit{Supplementary Table 1(a) continued.}}\\
\toprule
Score & Input representation & Ranking rule & Role \\
\midrule
\endhead
\bottomrule
\endlastfoot


Mean cosine
&
Control-referenced mean perturbation signatures
&
Descending cosine similarity between query and candidate mean-response vectors
&
Primary mean-signature baseline
\\

CMap-style cosine
&
Same control-referenced mean differential-expression vectors as mean cosine
&
Same cosine operation as the implemented mean-cosine baseline
&
Signature-retrieval reference; numerically identical to mean cosine in this implementation
\\

WTCS
&
Mean-collapsed differential-expression response
&
Rank-based weighted connectivity / enrichment score
&
Canonical rank-based signature baseline; related to but not algebraically identical to mean cosine
\\

Energy distance
&
Empirical control-referenced single-cell response populations
&
Candidates ranked by increasing energy distance to the query population
&
Primary population-level retrieval score
\\

RBF-MMD
&
Empirical response populations
&
Candidates ranked by increasing maximum mean discrepancy under an RBF kernel
&
Alternative distribution-sensitive score
\\

Sliced Wasserstein
&
Empirical response populations
&
Candidates ranked by increasing sliced Wasserstein distance using 128 random projections
&
Alternative distribution-sensitive score
\\

Subpopulation coverage
&
Response populations decomposed into response components
&
Component-level distances aggregated with temperature $\beta$; increasing $\beta$ gives progressively greater weight to the least well-matched response component
&
Population-sensitive score used for the partial-observation response-matching analyses
\\
\end{longtable}
\stabnote{The four score families become five population-level configurations because
subpopulation-coverage contributes two: coverage-mean is the $\beta\rightarrow0$ limit of the
aggregate and coverage-worst the $\beta\rightarrow\infty$ limit, both with energy distance as the
per-component base metric. The CMap-style cosine is numerically identical to mean cosine on
control-referenced mean vectors; its sign-flipped variant is scored in the partial-observation
experiments as an anti-correlation floor and is not reported here.}
\endgroup


\begingroup
\footnotesize
\begin{longtable}{@{}p{3.4cm}p{4.5cm}p{4.2cm}p{3.5cm}@{}}
\caption*{{\bfseries Supplementary Table 1 continued. (b) Retrieval task constructions.}}\\
\toprule
Task & Query and candidate construction & Candidate evaluation & Statistical unit \\
\midrule
\endfirsthead
\caption*{\textit{Supplementary Table 1(b) continued.}}\\
\toprule
Task & Query and candidate construction & Candidate evaluation & Statistical unit \\
\midrule
\endhead
\bottomrule
\endlastfoot


Observed-population retrieval over heterogeneous queries
&
A heterogeneous query population combines two response modes at weight $\alpha$; the library holds a held-out re-mixture at the same weight, the two single-mode populations, and 40 distractors, giving a 43-candidate library
&
Hit@1; only the both-mode candidate is relevant, so a single-mode population counts as a miss even when it carries one of the query's two response modes
&
Retrieval query; headline values macro-averaged over seven task--setting combinations
\\

Controlled heterogeneous queries
&
Two response modes are combined at mixing weight $\alpha$; controlled SciPlex3 analyses use HDAC- and JAK-associated response sources
&
Source identification and response-matching criteria
&
Query / task setting
\\

Partial-observation retrieval
&
Only the response information available under the controlled partial-observation protocol is exposed to the retriever; population and mean methods rank the identical candidate library once
&
The resulting fixed ranking is evaluated under response matching, MoA recovery and minority-state coverage
&
765 total queries
\\

Mechanism-evaluable subset
&
Same partial-observation rankings; restricted to queries for which the query mechanism remains represented among candidates
&
MoA-nDCG
&
600 queries; 165 excluded because MoA recovery is undefined
\\

Real-data minority-state projection
&
Top-ranked population and mean candidates evaluated on 239 task instances spanning constructed cross-line, within-line, Frangieh, CD34$+$ and predicted-population settings
&
Difference in minority-state coverage
&
Task-level difference; not interpreted as 239 independent biological replicates
\\

GDSC2-linked retrieval
&
Held-out SciPlex3 compound at 10\,$\mu$M used as query; remaining matched compounds form a 34--35 candidate library
&
Spearman agreement between RNA-derived ranking and external GDSC2 functional-similarity ordering
&
Per-query correlation summarized within SciPlex3 cell line
\\
\end{longtable}
\endgroup


\begingroup
\footnotesize
\begin{longtable}{@{}p{3.5cm}p{6.2cm}p{5.5cm}@{}}
\caption*{{\bfseries Supplementary Table 1 continued. (c) Evaluation criteria applied to fixed rankings.}}\\
\toprule
Criterion & Definition & Interpretation \\
\midrule
\endfirsthead
\caption*{\textit{Supplementary Table 1(c) continued.}}\\
\toprule
Criterion & Definition & Interpretation \\
\midrule
\endhead
\bottomrule
\endlastfoot


Hit@1
&
Indicator that the both-mode candidate, the held-out re-mixture at the query's own weight, is ranked first
&
Recovery of the query's response composition, not of its perturbation identity alone
\\

Response-matching regret
&
Energy-based welfare regret against the latent per-cell objective used to construct the controlled partial-observation task; reported as regret reduction relative to mean retrieval
&
Response fidelity; intentionally close to the population information used by the retriever
\\

MoA-nDCG
&
Normalized discounted cumulative gain rewarding candidates that share the query perturbation's annotated mechanism
&
Task-proximal biological criterion; undefined when the relevant mechanism is absent from the candidate library
\\

Minority-state coverage
&
Measures whether the selected candidate better preserves the response of the less-represented cellular state
&
Task-proximal measure of coverage of heterogeneous response structure
\\

GDSC2 functional similarity
&
Spearman correlation between cell-line-centred cross-cell-line AUC profiles of the query and candidate drugs
&
Externally measured functional similarity; not absolute potency
\\

Protein mean evaluator
&
Cosine similarity between control-referenced mean protein-response vectors
&
External evaluator preserving only mean protein response
\\

Protein population evaluator
&
Negative energy distance between control-referenced per-cell protein-response populations
&
External evaluator retaining cell-to-cell protein-response structure
\\

Response-magnitude matching
&
$-\left|\|\bar d_q\|-\|\bar d_c\|\right|$
&
Query-dependent scalar control that uses no within-population information
\\
\end{longtable}
\stabnote{Query states $s$ are obtained by $k$-means with $k=2$ on the query cells, the minority state
being the smaller cluster. A candidate $d$ with response population $P_d$ receives the welfare
$U(d)=\operatorname{agg}_s\big(-\mathcal{E}(P_d,Q_s)\big)$, where $\mathcal{E}$ is the energy
distance and the aggregator is the minimum over states in the worst-case form and the mean in the
mean form. The response-matching regret of a selected candidate $\hat d$ is
$\mathrm{Regret}(\hat d)=\max_{d\in\mathcal{L}}U(d)-U(\hat d)\ge 0$, taken over the observed library
$\mathcal{L}$, so the reference is the best available candidate rather than an unobservable
optimum. Minority-state coverage is
$\tfrac{1}{2}\big(1+\cos(\bar P_d,\bar Q_{\mathrm{min}})\big)$, the cosine between the candidate's
mean response and the mean response of the query's minority state rescaled to $[0,1]$. The same
definitions are given in the main Methods.}
\endgroup




% SUPPLEMENTARY TABLE 2
\begingroup
\footnotesize
\begin{longtable}{@{}p{2.7cm}p{2.5cm}p{3.0cm}p{3.4cm}p{4.0cm}@{}}
\caption*{{\large\bfseries Supplementary Table 2. Data resources, preprocessing and analytical roles.}}\\
\toprule
Resource & Scale & Perturbation/context structure & Representation and controls & Role in this study \\
\midrule
\endfirsthead
\caption*{\textit{Supplementary Table 2 continued.}}\\
\toprule
Resource & Scale & Perturbation/context structure & Representation and controls & Role in this study \\
\midrule
\endhead
\bottomrule
\endlastfoot


SciPlex3
&
276{,}325 cells; 2{,}000 genes
&
188 compounds shared across A549, K562 and MCF7
&
Log-transformed 2{,}000-gene representation; perturbation responses referenced to matched controls
&
Primary observed-population drug-retrieval setting; controlled within- and cross-context constructions; source of the GDSC2-linked transcriptomic rankings
\\

Frangieh Perturb-CITE-seq
&
218{,}023 RNA-profiled cells for retrieval and 218{,}331 barcode-matched cells for the RNA--protein evaluation; 2{,}000 genes; 20 retained surface proteins
&
239 CRISPR knockouts across Control, IFN$\gamma$ and Co-culture conditions
&
RNA and protein modalities joined by cell barcode; responses referenced to condition-matched controls
&
Independent genetic-perturbation retrieval setting and matched RNA--protein evaluator-shape experiment; not treated as a drug-retrieval benchmark
\\

CD34$+$ HSPC
&
33{,}984 cells; 2{,}000 genes
&
36 compounds
&
Log-transformed expression; DMSO-referenced perturbation responses
&
Negative-control retrieval setting for testing whether population-level information consistently improves ranking
\\

ZhaoSims2021 glioblastoma
&
10 patients; 6 compounds
&
Acute-slice culture and surgical biopsy with co-existing malignant, myeloid and oligodendrocyte compartments
&
Counts normalized to $10^4$ per cell, log1p-transformed and reduced to 2{,}000 highly variable genes; state-specific responses referenced within patient and compartment
&
Natural-tissue analysis of differential response, recoverability and state-specific response ordering; no retrieval benchmark is constructed
\\

Tahoe-100M plate 3
&
4{,}158{,}278 cells; 2{,}304 genes
&
50 cell lines $\times$ 93 compounds at one concentration, with DMSO\_TF controls
&
Released 2{,}304-gene log1p representation used without refitting; state-specific responses referenced to matched state controls
&
Two roles: large-scale measurement of differential response, recoverability and state-ordering agreement under intrinsic within-cell-line heterogeneity, and the substrate for the source-to-target intervention retrieval benchmark (Figs.~5 and~6)
\\

GDSC2
&
286 compounds across 969 cell lines in the retained release
&
35 compounds matched to the SciPlex3 panel; A549, K562 and MCF7 excluded from external profiles
&
Fitted dose--response AUC; primary functional similarity uses cell-line-centred cross-line response profiles
&
External functional evaluation of fixed SciPlex3 RNA-derived rankings
\\

HIR-Bench
&
13{,}440 synthetic instances
&
672 parameter configurations $\times$ 20 seeds in a two-subpopulation latent-welfare model
&
Synthetic observed, predicted-mean and predicted-structure information conditions
&
Controlled demonstration that population structure can alter candidate preference; not interpreted as a biological generative model
\\
\end{longtable}
\stabnote{The resources above answer distinct questions and are not pooled as independent
replicates of a single benchmark. Dataset-specific cell-number thresholds, subsampling rules and
statistical units are given in the corresponding Methods sections.}
\endgroup




% SUPPLEMENTARY TABLE 3
\begingroup
\small
\begin{longtable}{@{}lcc@{}}
\caption*{{\large\bfseries Supplementary Table 3. Sensitivity of the GDSC2 functional-similarity evaluation to response-profile construction.}}\\
\caption*{The primary external functional criterion compares drug--drug response profiles after removing the shared cell-line sensitivity effect. For cell line $l$ and drug $d$, \[ \widetilde{\mathrm{AUC}}_{d,l} = \mathrm{AUC}_{d,l} - \frac{1}{|\mathcal D_l|} \sum_{d'\in\mathcal D_l} \mathrm{AUC}_{d',l}, \] where $\mathcal D_l$ contains all GDSC2 compounds measured in that cell line. Drug--drug functional similarity is then the Spearman correlation between centred cross-cell-line response profiles. The table compares this primary construction with the corresponding uncentred AUC-profile similarity.}\\
\toprule
RNA-derived ranking rule
&
\shortstack{Cell-line-centred\\functional similarity}
&
\shortstack{Uncentred\\functional similarity}
\\
\midrule
\endfirsthead
\caption*{\textit{Supplementary Table 3 continued.}}\\
\toprule
RNA-derived ranking rule
&
\shortstack{Cell-line-centred\\functional similarity}
&
\shortstack{Uncentred\\functional similarity}
\\
\midrule
\endhead
\bottomrule
\endlastfoot


%% REISSUED 2026-09-04 with the main manuscript's \ENERGYFUNC and \ENERGYUNCEN. These were
%% +0.276 and +0.0967, the V-statistic values; Fig. 3h had already been rebuilt and drew +0.265.
%% Only the energy row moves: every other row of this table ranks by a mean vector or a scalar
%% and never touches the energy kernel.
Energy distance
&
\textbf{+0.265}
&
+0.0902
\\

Mean cosine, control referenced
&
+0.083
&
\textbf{+0.138}
\\

Mean cosine, treated mean only
&
+0.241
&
+0.092
\\

Response-magnitude matching
&
+0.232
&
+0.217
\\

Candidate magnitude only
&
$-$0.330
&
+0.135
\\

Potency matching diagnostic
&
+0.399
&
+0.290
\\

\midrule

Energy, magnitude-adjusted partial rank correlation
&
+0.105
&
n/a
\\
\end{longtable}
\stabnote{Values are median per-query Spearman correlations over 103 leave-one-drug-out queries across the
three SciPlex3 cell lines. The primary centred functional criterion is constructed from GDSC2 AUC
profiles after subtracting each cell line's mean response across all 286 compounds in the retained
release. Drug pairs sharing fewer than 100 cell lines are excluded; 591 usable drug pairs remain,
with a median of 909 shared cell lines per pair.

\smallskip\noindent
The change between centred and uncentred constructions shows that external evaluation is itself
sensitive to how the functional measurement is summarized. Under the centred criterion, energy
retrieval exceeds the control-referenced mean-signature baseline, but much of this association is
also reproduced by response-magnitude matching. Under the uncentred criterion, the ordering between
energy and the mean-signature baseline reverses. Neither construction is interpreted as a direct
measure of therapeutic utility.}
\endgroup







% SUPPLEMENTARY TABLE 4
\begingroup
\footnotesize
\begin{longtable}{@{}p{2.5cm}p{3.3cm}p{3.2cm}p{3.6cm}p{3.1cm}@{}}
\caption*{{\large\bfseries Supplementary Table 4. Perturbation-predictor implementations and their role in the differential-response analysis.}}\\
\toprule
Predictor & Model form & Predicted-population construction & Use in this study & Interpretation / limitation \\
\midrule
\endfirsthead
\caption*{\textit{Supplementary Table 4 continued.}}\\
\toprule
Predictor & Model form & Predicted-population construction & Use in this study & Interpretation / limitation \\
\midrule
\endhead
\bottomrule
\endlastfoot


Average-effect
&
Per-drug mean perturbation delta
&
Predicted delta applied to the real control cells of the query context
&
Predict-then-rank candidate source and within-context differential-response reference
&
Additive by construction; state-specific induced-response cosine equals 1
\\

Nearest-neighbour
&
Transfers the perturbation effect from the training context with the most similar untreated control state
&
Transferred effect applied to real query-context control cells
&
Predict-then-rank candidate source
&
Tests simple context transfer rather than a general virtual-cell architecture
\\

Linear-latent
&
Additive latent arithmetic in an encoded representation followed by linear decoding
&
Observed control cells encoded, shifted in latent space and decoded individually
&
Predict-then-rank candidate source and additive latent reference
&
Can preserve baseline cellular heterogeneity while remaining additive in perturbation effect
\\

scGen
&
Variational-autoencoder perturbation model with latent arithmetic
&
Evaluated on the same within-context state partition and state-specific response definition used for other predictors
&
Additional differential-response diagnostic reported in the main text
&
Mean induced-response cosine 0.972 over 30 cell-line and drug combinations
\\

Entropic optimal-transport map
&
Non-parametric Sinkhorn mapping that permits cell-dependent displacements
&
Control cells mapped individually to predicted treated states
&
Additional non-additive differential-response diagnostic; also tested as a candidate-response source
&
Within-context induced-response cosine approximately $0.900\pm0.107$ against the map's own reconstruction baseline and $0.380$ against the real control; cross-context candidate predictions were too inaccurate for a clean retrieval test
\\
\end{longtable}
\endgroup


%% THE FIRST THREE ROWS WERE A CHECKLIST, NOT A SPECIFICATION, until 2026-09-04. They named the
%% fields a reproducible entry has to carry ("drug-effect estimation set, context handling, cell
%% cap, and exact construction") without filling any of them in, while the scGen and
%% optimal-transport rows below were already written out. They are now filled from the code that
%% produced the reported runs: src/baselines/average_effect_predictor.py,
%% src/baselines/nearest_neighbor_predictor.py, the _CPALinearLatent class in
%% src/baselines/scgen_predictor.py (recorded in the run provenance as backend cpa_linear), and
%% the runner src/experiments/exp09_predict_then_rank.py.
\begingroup
\footnotesize
\begin{longtable}{@{}p{3.0cm}p{11.8cm}@{}}
\caption*{{\bfseries Supplementary Table 4 continued. Implementation of each predictor.}}\\
\toprule
Predictor & Implementation \\
\midrule
\endfirsthead
\caption*{\textit{Supplementary Table 4 continued.}}\\
\toprule
Predictor & Implementation \\
\midrule
\endhead
\bottomrule
\endlastfoot


Average-effect
&
Per drug $d$, $\widehat{\delta}_d=\operatorname{mean}_{c}\big(\operatorname{mean}(\text{treated}_{c,d})-\operatorname{mean}(\text{control}_{c})\big)$
over the SciPlex3 cell lines $c$, a line contributing only where it has at least two treated cells
for $d$; one signature per drug, independent of the query context. The predicted population is that
single signature added to each of the query context's real control cells, so every cell is shifted
identically. No reconstruction step, so the model's own zero-effect baseline is the control cells
themselves. A drug absent from the fit is scored at $-10^{6}$ under every rule rather than imputed
\\

Nearest-neighbour
&
$c^{*}=\arg\min_{c}\big(1-\cos(\bar{x}_{\text{control}}^{\text{query}},\bar{x}_{\text{control}}^{c})\big)$
over the training lines, then $\widehat{\delta}_d=\operatorname{mean}(\text{treated}_{c^{*},d})-\operatorname{mean}(\text{control}_{c^{*}})$.
Distances are compared strictly, so an exact tie is resolved by context iteration order; where
$d$ was not measured in $c^{*}$ the signature falls back to the mean of $d$'s signatures over the
remaining admissible lines. Population synthesis as for average-effect. The two query-composing cell lines remained eligible donor
contexts, so this experiment is interpreted as an in-sample ceiling rather than as a transfer
evaluation
\\

Linear-latent
&
PCA with 30 components fitted, unsupervised and once, on up to 20{,}000 cells drawn from the full
SciPlex3 matrix at seed 0. Each line's control latent mean is stored; each drug's latent delta is
the mean over lines of (mean latent of treated $-$ that line's control latent mean), again
requiring at least two treated cells. Prediction encodes each real control cell of the query
context, adds the drug's latent delta to that cell's own latent code, and decodes each cell by the
PCA inverse transform. The decoder is linear, so the composite map is
$\mathrm{dec}(\mathrm{enc}(x)+\mathrm{d}z)=\mathrm{P}x+\mathrm{P}\,\mathrm{d}z$ and the predictor is
additive in gene space and deterministic. Because it reconstructs, its own zero-effect baseline is
the PCA reconstruction of the control cells rather than those cells
\\

scGen
&
Published \texttt{scgen.SCGEN} on \texttt{scvi-tools}; trained on the SciPlex3 tensor, batch key
= perturbation, label key = cell line; 30 latent dimensions; batch size 512; at most 40 epochs with
early stopping after 15 epochs without improvement; fixed seed; per-perturbation latent delta added
to the encoded control mean and decoded
\\

Optimal transport
&
Entropic (Sinkhorn) coupling in a 30-component PCA latent; cost rescaled by its maximum; uniform
marginals; regularization $0.05$; at most 500 iterations; fitted per non-query context centred on
that context's own control mean, with the query context excluded; barycentric projection extended to
the query's control cells by nearest-source matching and decoded by the PCA inverse transform
\\
\end{longtable}
\stabnote{Predictors are used as candidate-response sources or information diagnostics rather than as a
benchmark of perturbation-prediction architectures. The reported predict-then-rank experiment uses
predictors fitted on the full SciPlex3 matrix and is therefore interpreted as an in-sample ceiling,
not as a cross-context generalization benchmark.

\smallskip\noindent
Constants shared by every row of the predict-then-rank run: 720 cross-line queries, formed from the
three cell-line pairs (K562--A549, A549--MCF7, K562--MCF7) at mixing weights $\alpha\in\{0.5,0.7,0.9\}$
over 10 divergent drugs and 8 seeds; 200 query cells per query; a 31-drug candidate library holding
the true drug and 30 other shared drugs; 200 predicted cells synthesized per candidate; query
construction seeded at $1000+s$ and distractor selection at $7000+s$ for $s=0,\dots,7$. The
within-context differential-response measurements in Supplementary Note~4 are a separate run over
30 cell-line and drug combinations with 400 control cells per context split into two states, a
state requiring 60 cells and a treated arm 40, at seed 0. Nearest-neighbour appears only in the
predict-then-rank run and scGen only in the differential-response run; the entropic
optimal-transport map appears in both.}
\endgroup

%% SUPPLEMENTARY REFERENCES, added 2026-09-06. The SI loaded natbib from the start but had no
%% \bibliography, so the methods named only here, Leiden, HDBSCAN, the entropic map and its POT
%% implementation, and scvi-tools, could not be credited anywhere. This list numbers from 1 and is
%% independent of the main reference list, which is the normal arrangement for a separate SI file:
%% a work cited in both documents carries a different number in each.
\clearpage
\renewcommand{\refname}{Supplementary References}
\bibliographystyle{unsrtnat}
\bibliography{references}

\end{document}
